\documentclass[10pt,a4paper]{amsart}
\usepackage{amsmath,amssymb,amsthm,graphicx}
% Retained minimal preamble for the Stacks Project chapter more-algebra.tex (More on Algebra).
% The chapter file pinned-original-source.tex carries only \input{preamble} on its first line, so no
% preamble text lives inside the chapter itself.  The Stacks Project uses one single shared preamble for
% every chapter; the snapshot stages it identically as sources/stacks-chow/preamble.tex and
% sources/stacks-etale-cohomology/preamble.tex.  The block below is copied mechanically, line for line,
% from sources/stacks-chow/preamble.tex lines 14-18 (the Xy-pic setup) and 154-217 (the native theorem
% declarations together with the mathematical macros).  Nothing was added, removed or redefined; the
% wrapper class and the amsmath/amssymb/amsthm/graphicx base packages are supplied by the assembler.
% Chapter: more-algebra.tex -- More on Algebra.
% For commutative diagrams we use Xy-pic
\usepackage[all]{xy}

% We use 2cell for 2-commutative diagrams.
\xyoption{2cell}

% Theorem environments.
%
\theoremstyle{plain}
\newtheorem{theorem}[subsection]{Theorem}
\newtheorem{proposition}[subsection]{Proposition}
\newtheorem{lemma}[subsection]{Lemma}

\theoremstyle{definition}
\newtheorem{definition}[subsection]{Definition}
\newtheorem{example}[subsection]{Example}
\newtheorem{exercise}[subsection]{Exercise}
\newtheorem{situation}[subsection]{Situation}

\theoremstyle{remark}
\newtheorem{remark}[subsection]{Remark}
\newtheorem{remarks}[subsection]{Remarks}

\numberwithin{equation}{subsection}

% Macros
%
\def\lim{\mathop{\mathrm{lim}}\nolimits}
\def\colim{\mathop{\mathrm{colim}}\nolimits}
\def\Spec{\mathop{\mathrm{Spec}}}
\def\Hom{\mathop{\mathrm{Hom}}\nolimits}
\def\Ext{\mathop{\mathrm{Ext}}\nolimits}
\def\SheafHom{\mathop{\mathcal{H}\!\mathit{om}}\nolimits}
\def\SheafExt{\mathop{\mathcal{E}\!\mathit{xt}}\nolimits}
\def\Sch{\mathit{Sch}}
\def\Mor{\mathop{\mathrm{Mor}}\nolimits}
\def\Ob{\mathop{\mathrm{Ob}}\nolimits}
\def\Sh{\mathop{\mathit{Sh}}\nolimits}
\def\NL{\mathop{N\!L}\nolimits}
\def\CH{\mathop{\mathrm{CH}}\nolimits}
\def\proetale{{pro\text{-}\acute{e}tale}}
\def\etale{{\acute{e}tale}}
\def\QCoh{\mathit{QCoh}}
\def\Ker{\mathop{\mathrm{Ker}}}
\def\Im{\mathop{\mathrm{Im}}}
\def\Coker{\mathop{\mathrm{Coker}}}
\def\Coim{\mathop{\mathrm{Coim}}}

% Boxtimes
%
\DeclareMathSymbol{\boxtimes}{\mathbin}{AMSa}{"02}

%
% Macros for moduli stacks/spaces
%
\def\QCohstack{\mathcal{QC}\!\mathit{oh}}
\def\Cohstack{\mathcal{C}\!\mathit{oh}}
\def\Spacesstack{\mathcal{S}\!\mathit{paces}}
\def\Quotfunctor{\mathrm{Quot}}
\def\Hilbfunctor{\mathrm{Hilb}}
\def\Curvesstack{\mathcal{C}\!\mathit{urves}}
\def\Polarizedstack{\mathcal{P}\!\mathit{olarized}}
\def\Complexesstack{\mathcal{C}\!\mathit{omplexes}}
% \Pic is the operator that assigns to X its picard group, usage \Pic(X)
% \Picardstack_{X/B} denotes the Picard stack of X over B
% \Picardfunctor_{X/B} denotes the Picard functor of X over B
\def\Pic{\mathop{\mathrm{Pic}}\nolimits}
\def\Picardstack{\mathcal{P}\!\mathit{ic}}
\def\Picardfunctor{\mathrm{Pic}}
\def\Deformationcategory{\mathcal{D}\!\mathit{ef}}
\usepackage{cleveref}
\crefname{originallemma}{original lemma}{original lemmas}
\Crefname{originallemma}{Original lemma}{Original lemmas}
\makeatletter
\newlabel{lemma-better-cut-complex-in-two}{{19744}{source}{}{sourcepointer-1}{}}
\newlabel{lemma-better-cut-complex-in-two@cref}{{[originallemma][0][]19744}{[1][0][]source}{}{}{}}
\makeatother
\title{A vanishing cohomology group of a pseudo-coherent complex at a prime cuts the localized complex into an upper and a lower part}
\author{The Stacks Project Authors}
\date{stacks-project:more-algebra.tex:local-snapshot-20261012; GNU FDL 1.2 or later}
\begin{document}
\maketitle
\setcounter{subsection}{1}
\noindent\textit{Editorial note.} All delivered mathematical text is the author's own, copied from the pinned Stacks Project chapter file more-algebra.tex; no statement or proof step was rewritten, reordered or removed.\par\medskip
\noindent\textit{Editorial note.} The chapter file pinned-original-source.tex carries only \textbackslash{}input\{preamble\} on its first line, so no preamble text lives inside the chapter. The retained file sources/stacks-more-algebra/retained-preamble.tex is a mechanical, line-for-line copy of the project's single shared preamble as staged in the snapshot; nothing was added, removed or redefined and no package was added.\par\medskip
\noindent\textit{Editorial note.} The source URL is the source-level Stacks Project URL, as required when more than one card is produced from one source; the per-claim tag is carried by claim\_key.\par\medskip
\noindent\textit{Editorial note.} The author opens the proof by naming a special case of another lemma and then says 'A direct proof is as follows.', supplying that proof in full inside the counted span; the single \textbackslash{}ref pointer to lemma-better-cut-complex-in-two is recorded as a non-load-bearing navigation pointer, and the direct argument that replaces it stays inside the counted proof span.\par\medskip
\noindent\textit{Allowed standard prerequisites (editorial).}\begin{itemize}
\item Pseudo-coherent complexes and their representation by complexes of finite free modules. the hypothesis at statement lines 19869-19870 and used at proof lines 19887-19888. Legitimises the reduction to a bounded above complex of finite free modules on which the matrix argument runs.
\item Truncations tau\_\{>= n\} and tau\_\{<= n\} of a complex and tor-amplitude. the conclusion at statement lines 19875-19879. Names the two pieces of the asserted decomposition and the property of the upper one.
\item Rank of a matrix over a field and nonvanishing of a minor. used at proof lines 19896-19899. Converts the rank statement modulo p into the two minors a and b whose product f is the localizing element.
\end{itemize}\medskip
\ifdefined\hypertarget\hypertarget{sourcepointer-1}{}\fi
\noindent\textit{Source reference.} Original lemma 19744: sources/stacks-more-algebra/pinned-original-source.tex, source lines 19744--19744. Same-chapter navigation pointer; the statement it names is standard for the section and the argument that consumes it stays inside the counted proof span.\par\medskip
\begin{lemma}
\label{lemma-cut-complex-in-two}
Let $R$ be a ring. Let $\mathfrak p \subset R$ be a prime ideal.
Let $i \in \mathbf{Z}$. Let $K^\bullet$ be a pseudo-coherent complex
of $R$-modules such that
$H^i(K^\bullet \otimes_R^{\mathbf{L}} \kappa(\mathfrak p)) = 0$.
Then there exists an $f \in R$, $f \not \in \mathfrak p$
and a canonical direct sum decomposition
$$
K^\bullet \otimes_R R_f =
\tau_{\geq i + 1}(K^\bullet \otimes_R R_f) \oplus
\tau_{\leq i - 1}(K^\bullet \otimes_R R_f)
$$
in $D(R_f)$ with $\tau_{\geq i + 1}(K^\bullet \otimes_R R_f)$ a perfect
complex with tor-amplitude in $[i + 1, \infty]$.
\end{lemma}

\begin{proof}
This is an often used special case of
Lemma \ref{lemma-better-cut-complex-in-two}.
A direct proof is as follows.
We may assume that $K^\bullet$ is a bounded above
complex of finite free $R$-modules. Let us inspect what is happening
in degree $i$:
$$
\ldots \to K^{i - 2} \to R^{\oplus l}
\to R^{\oplus m} \to R^{\oplus n} \to K^{i + 2} \to \ldots
$$
Let $A$ be the $m \times l$ matrix corresponding to $K^{i - 1} \to K^i$
and let $B$ be the $n \times m$ matrix corresponding to $K^i \to K^{i + 1}$.
The assumption is that $A \bmod \mathfrak p$ has rank $r$ and that
$B \bmod \mathfrak p$ has rank $m - r$. In other words, there is some
$r \times r$ minor $a$ of $A$ which is not in $\mathfrak p$ and there is
some $(m - r) \times (m - r)$-minor $b$ of $B$ which is not in $\mathfrak p$.
Set $f = ab$. Then after inverting $f$ we can find direct sum decompositions
$K^{i - 1} = R^{\oplus l - r} \oplus R^{\oplus r}$,
$K^i = R^{\oplus r} \oplus R^{\oplus m - r}$,
$K^{i + 1} = R^{\oplus m - r} \oplus R^{\oplus n - m + r}$
such that the module map $K^{i - 1} \to K^i$ kills of
$R^{\oplus l - r}$ and induces an isomorphism of $R^{\oplus r}$ onto the
corresponding summand of $K^i$ and such that the module map $K^i \to K^{i + 1}$
kills of $R^{\oplus r}$ and induces an isomorphism of $R^{\oplus m - r}$
onto the corresponding summand of $K^{i + 1}$. Thus $K^\bullet$ becomes
quasi-isomorphic to
$$
\ldots \to K^{i - 2} \to R^{\oplus l - r}
\to 0 \to R^{\oplus n - m + r} \to K^{i + 2} \to \ldots
$$
and everything is clear.
\end{proof}

\end{document}
